\noindent
Discrete feature growth (DFG) is one greedy loop: fit the labels by least squares on the features found so far, and admit every Fourier character or signed-sum indicator whose correlation with the residual passes a union-bound test over its class. DFG v1.3, run alone here, is the one new learning algorithm in v3, the stacked portfolio of this study's previous part, and the source of its Boolean results.

On its held-out suite, 12 fresh Boolean targets, DFG needs clearly less data than the per-target best of three network variants on 10 targets and more on none: the mean $\log_2$ ratio of $\nstar$, the training-set size at 10\% test error, is $-1.29$. Over 126 in-sample targets it is $-1.87$. DFG's $\nstar$ grows with the input dimension $d$ with exponents 0.01--0.50, the networks' with 0.91--2.40. The low-degree algorithm, with DFG's characters and test but no search, needs $3.1\times$ more data and is level with the networks.

DFG's models are short lists of named features, often the target's exact structure. On parities a proved criterion predicts every fit with $n\ge32$; on signed-sum targets the search, not the test, sets the data cost, and residue classes mod 3, with almost no weight at Fourier degrees 1--2, are never found. DFG admits characters in order of detectability, the networks in order of construction. A post-hoc core with five parts nearly matches DFG at a tenth of its FLOPs. DFG needs $\pm1$ inputs: real-valued data, where earlier adaptations lost to a regularized MLP, is not covered.
